\documentclass[10pt, a4paper]{article}

% --- UNIVERSAL PREAMBLE BLOCK ---
\usepackage[a4paper, landscape, top=1.5cm, bottom=1.5cm, left=1.2cm, right=1.2cm]{geometry}
\usepackage{fontspec}

% English setup
\usepackage[english]{babel}

% Set default font to Noto Sans
\babelfont{rm}{Noto Sans}

% --- REQUIRED PACKAGES FOR TABLE FORMATTING ---
\usepackage{longtable}
\usepackage{booktabs}
\usepackage{array}
\usepackage{hyperref}

\begin{document}

\begin{center}
\scriptsize
\renewcommand{\arraystretch}{1.5}
\begin{longtable}{
  >{\raggedright\arraybackslash\bfseries}p{2.8cm}
  >{\raggedright\arraybackslash}p{2.8cm}
  >{\raggedright\arraybackslash}p{1.8cm}
  >{\raggedright\arraybackslash}p{2.5cm}
  >{\raggedright\arraybackslash}p{6.8cm}
  >{\raggedright\arraybackslash}p{6.0cm}
}
\caption{The Exhaustive Taxonomic Deconstruction of Combinatorial Optimization Operators} \label{tab:operator_taxonomy_ultimate} \\
\toprule
\normalfont\textbf{Operator} & \normalfont\textbf{Origin / Key Ref.} & \normalfont\textbf{Category} & \normalfont\textbf{Primary Mechanism} & \normalfont\textbf{Topological Effect / Structural Description} & \normalfont\textbf{Computational Nuance} \\
\midrule
\endfirsthead

\multicolumn{6}{c}%
{{\bfseries \tablename\ \thetable{} -- continued from previous page}} \\
\toprule
\normalfont\textbf{Operator} & \normalfont\textbf{Origin / Key Ref.} & \normalfont\textbf{Category} & \normalfont\textbf{Primary Mechanism} & \normalfont\textbf{Topological Effect / Structural Description} & \normalfont\textbf{Computational Nuance} \\
\midrule
\endhead

\midrule \multicolumn{6}{r}{{Continued on next page}} \\
\endfoot

\bottomrule
\endlastfoot

% --- CROSSOVER & POPULATION OPERATORS ---
{Ordered (OX)} & Davis (1985) & Crossover & Substring Copy & Copies a contiguous substring from one parent to preserve sequence; remaining loci filled by second parent. & $\mathcal{O}(n)$. Fast, but susceptible to premature convergence vs. graph-based operators. \\ \midrule
{Partially Mapped (PMX)} & Goldberg \& Lingle (1985) & Crossover & Mapping & Preserves absolute positions from one parent and uses a mapping array to resolve conflicts from the other. & $\mathcal{O}(n)$. Essential for strict permutation problems (TSP) to avoid duplicate vertices. \\ \midrule
{Cycle (CX)} & Oliver et al. (1987) & Crossover & Cyclic Preserv. & Identifies cycles between parents to ensure every node's position is inherited from exactly one parent. & $\mathcal{O}(n)$. Highly disruptive to edges but perfectly preserves absolute positional indices. \\ \midrule
{Position Indep. (PIX)} & Syswerda (1991) & Crossover & Binary Masking & Disregards absolute spatial sequencing; uses a randomized binary mask to pull vertices from either parent. & $\mathcal{O}(n)$. Creates massive permutation violations; strictly requires robust external repair. \\ \midrule
{Edge Recomb. (ERX)} & Whitley et al. (1989) & Crossover & Adj. Matrix & Prioritizes adjacent edge preservation over positional mapping via a unified neighborhood matrix. & $\mathcal{O}(n^2)$. High overhead due to dynamic matrix updates; scales poorly for massive VRPs. \\ \midrule
{Gen. Partition (GPX)} & Whitley et al. (1998) & Crossover & Graph Tunneling & Overlays parent solutions into a union multigraph to identify and evaluate independent recombining partitions. & $\mathcal{O}(n)$. Tunnels directly between optima by exploiting mathematically independent subgraphs. \\ \midrule
{Edge Assembly (EAX)} & Nagata \& Kobayashi (1997) & Crossover & Sub-tour Patching & Merges parent edges into AB-cycles, generating intermediate disjoint sub-tours patched via spanning trees. & $\mathcal{O}(n \log n)$. Extremely powerful crossover in existence for TSP and symmetric routing. \\ \midrule
{Sel. Route Ex. (SREX)} & Homberger \& Gehring (1999) & Crossover & Macro Exchange & Evaluates holistic constraints by replacing an intact subset of vehicle routes with overlapping parent routes. & $\mathcal{O}(n \log n)$. Essential for PDPTW/VRP; preserves fragile interdependent vehicle systems. \\ \midrule
{Path Relinking} & Glover (1997) & Population & Trajectory Tunnel & Iteratively forces a source solution's topological attributes to systematically match a guiding elite solution. & $\mathcal{O}(n^2)$ per step. Bridges genetic recombination directly with deterministic trajectory search. \\ \midrule

% --- DESTROY (RUIN) OPERATORS ---
{Random Removal} & Ropke \& Pisinger (2006) & Destroy & Stoch. Excision & Uniformly and randomly removes a predetermined fraction of vertices from the current routing deployment. & Baseline unguided diversification; mathematically simple but breaks deterministic looping. \\ \midrule
{Worst Removal} & Ropke \& Pisinger (2006) & Destroy & Cost Excision & Targets and extracts vertices that contribute the most disproportionately high cost to the objective function. & Eliminates accumulated algorithmic inefficiencies; highly effective when used iteratively. \\ \midrule
{Cluster Removal} & Smilowitz et al. (2004) & Destroy & Spatial Partit. & Utilizes algorithms (e.g., k-means, Kruskal's) to simultaneously excise entire geographic/spatial clusters. & Forces complete regional topological redesigns rather than minor adjustments. \\ \midrule
{Shaw (Relatedness)} & Shaw (1997/1998) & Destroy & Affinity Excision & Removes sets of highly interconnected nodes based on a composite of distance, time-windows, and load. & Exploits the premise that swapping related customers yields highly optimized re-insertions. \\ \midrule
{Time-Window Rem.} & Ropke \& Pisinger (2006) & Destroy & Temporal Targ. & Excises nodes whose scheduled service times aggressively tightly cluster, causing temporal bottlenecking. & Crucial for VRPTW; frees up timeline slack to allow downstream shifting. \\ \midrule
{String Removal} & Christiaens et al. (2020) & Destroy & Seq. Excision & Removes contiguous sequential strings of nodes from a given route. Core to SISR framework. & Preserves overarching framework while creating localized internal voids for deep repair. \\ \midrule
{Route Removal} & Ropke \& Pisinger (2006) & Destroy & Macro Emptying & Empties entire specific vehicles, forcing all associated customers back into the unassigned node pool. & Critical algorithmic driver for fleet size minimization and heavy demand consolidation. \\ \midrule
{Neighbor (Graph)} & Ropke \& Pisinger (2006) & Destroy & Adj. Excision & Targets and removes vertices directly, sequentially adjacent to a randomly selected node in the tour. & Disrupts localized sequential chains trapped in sub-optimal graph permutations. \\ \midrule
{Hist. Knowledge} & Demir et al. (2012) & Destroy & Adaptive Mem. & Tracks algorithmic history to deliberately target edges that frequently lead to sub-optimal local minima. & Leverages long-term systemic failures to intelligently guide topological destruction. \\ \midrule
{Sector (Sweep)} & Gillett \& Miller (1974) & Destroy & Geom. Sweeping & Removes all nodes and routes intersecting a randomly chosen 2D angular sweep from the depot. & Forces aggressive physical geographic realignments of intersecting routes. \\ \midrule
{Type I Unstringing} & Gendreau et al. (1992) & Destroy & Clean Extract & Removes node $V_k$ strictly between $V_i, V_j$ by creating edge $(V_i, V_j)$, assuming no overlapping endpoints. & Formalized topological unlinking ensuring residual graphs maintain strictly feasible paths. \\ \midrule
{Type II Unstringing} & Gendreau et al. (1992) & Destroy & Complex Extract & Breaks multiple non-adjacent edges simultaneously to extract a node without isolating sub-tours. & Symmetric unstringing that mathematically guarantees no isolated loops are created. \\ \midrule
{Type III Unstringing} & Gendreau et al. (1992) & Destroy & Directed Extract & Asymmetric inverse of Type I; reverses extraction logic to explicitly preserve directional graph flow. & Mandatory for environments where travel matrices are asymmetrical. \\ \midrule
{Type IV Unstringing} & Gendreau et al. (1992) & Destroy & Complex Asym. & Executes complex extractions under strict directional constraints while preventing internal circuit creation. & The asymmetric equivalent of Type II, preserving strict connectivity for directed VRPs. \\ \midrule

% --- REPAIR OPERATORS ---
{Greedy Insertion} & Clarke \& Wright (1964) & Repair & Immed. Minim. & Iteratively inserts nodes resulting in the absolute smallest immediate increase in objective cost. & Fast but fatally short-sighted; strands difficult outlier nodes into inefficient late positions. \\ \midrule
{2-Regret Insertion} & Potvin \& Rousseau (1993) & Repair & Penalty Diff. & Inserts based on 'regret': the cost difference between a node's absolute best and 2nd-best feasible spots. & Embeds critical look-ahead capability; prioritizes nodes facing severe exponential penalties. \\ \midrule
{$k$-Regret Insertion} & Ropke \& Pisinger (2006) & Repair & Deep Regret & Generalizes regret to calculate the differential between the best and the $k^{th}$ best available position. & Broadens predictive scope at the cost of heavily repeated, localized insertion evaluations. \\ \midrule
{Greedy w/ Blinks} & Christiaens et al. (2020) & Repair & Stoch. Accept. & Periodically "blinks" (probabilistically ignores) the mathematical optimal insertion point. & Simulated Annealing logic injected into repair to prevent looping and maintain diversity. \\ \midrule
{Nearest Neighbor} & Rosenkrantz et al. (1977) & Repair & Prox. Linking & Appends the geographically nearest unrouted customer directly to the end of the active partial tour. & Extremely fast $\mathcal{O}(n^2)$; inherently flawed for dense clustering due to "crossing-path" artifacts. \\ \midrule
{Savings-based} & Clarke \& Wright (1964) & Repair & Struct. Savings & Prioritizes insertions maximizing geographic linking savings compared to independent back-and-forth trips. & Foundational OR heuristic; highly effective when paired with Cluster Removal. \\ \midrule
{Deep (Look-ahead)} & Iori et al. (2007) & Repair & Cascading Eval. & Evaluates multi-tiered downstream effects of insertions (e.g., synchronized temporal horizons). & Computationally exhaustive; strictly required for interdependent constraints (e.g., synchronization). \\ \midrule
{Type I Stringing} & Gendreau et al. (1992) & Repair & Simple Stringing & Safely injects node $V_x$ between $V_i, V_j$ by executing straightforward edge reconnections. & Direct inverse formulation to the Type I unstringing heuristic framework. \\ \midrule
{Type II Stringing} & Gendreau et al. (1992) & Repair & Complex Reconn. & Re-maps multiple fractured edges simultaneously when injecting a node to prevent isolated loop generation. & Ensures structurally sound repair in multi-edge manipulation scenarios. \\ \midrule
{Type III Stringing} & Gendreau et al. (1992) & Repair & Asym. String I & Maps to the graph theory operations of Type I but specifically calibrated for strict directional edge reconnections. & Accommodates asymmetric CVRP matrix environments. \\ \midrule
{Type IV Stringing} & Gendreau et al. (1992) & Repair & Asym. String II & Mathematically re-strings complex, multi-fractured directed graphs without violating asymmetric constraints. & Advanced, structurally-aware graph re-assimilation. \\ \midrule


% --- INTRA-ROUTE OPERATORS (LOCAL SEARCH) ---
{Relocate (1-0)} & Savelsbergh (1992) & Intra-route & Position Shift & Extracts a single node and reinserts it into a different location sequentially within the exact same route. & $\mathcal{O}(n^2)$ neighborhood. Fundamental, high-speed foundational step for trajectory descent. \\ \midrule
{Swap (1-1)} & Savelsbergh (1992) & Intra-route & Pair Exchange & Directly exchanges the spatial locations of two distinct customer nodes within the same local sequence. & $\mathcal{O}(n^2)$ neighborhood. Instantaneous structural differential evaluation. \\ \midrule
{2-opt} & Croes (1958) & Intra-route & Sub-tour Rev. & Deletes two non-adjacent intra-route edges and reconnects, reversing the orientation of the intermediary sequence. & $\mathcal{O}(n^2)$. Primary mechanism utilized for rapidly untangling geometrically crossed paths. \\ \midrule
{3-opt} & Lin (1965) & Intra-route & Complex Rev. & Removes three edges and explores up to eight geometric reconnection permutations. & Escalates complexity geometrically to resolve tightly tangled local pathings. \\ \midrule
{$k$-opt} & Lin (1965) & Intra-route & Edge Shift & Expands structural edge exchange to $k$ breaking points across the route graph. & PSPACE-hard theoretical complexity; application gated by factorial scaling bounds. \\ \midrule
{Lin-Kernighan (LK)} & Lin \& Kernighan (1973) & Intra-route & Dynamic $k$-opt & Adaptive variable-depth search chaining alternating edge swaps to build complex $k$-opt moves. & The most powerful TSP heuristic. Prunes space via minimum spanning trees for tractability. \\ \midrule
{Or-opt} & Or (1976) & Intra-route & Chain Reloc. & Isolates intact sub-chains of nodes (length 1-3) and repositions the unbroken chain. & Preserves pre-optimized internal linkage efficiencies while enhancing macroscopic position. \\ \midrule
{GENI} & Gendreau et al. (1992) & Repair & Assimilation & Models node insertion strictly between nearest neighbors, triggering immediate localized re-optimization. & Approximates deep structural improvements of complex maneuvers bypassing combinatorics. \\ \midrule
{$k$-Permutation} & Christofides \& Eilon (1969)& Intra-route & Brute-force & Isolates a precise string of $k$ nodes and strictly evaluates all $k!$ sequential permutations. & Guarantees localized optimality but completely gated by factorial scaling ($k \le 5$). \\ \midrule
{3-Permutation} & Christofides \& Eilon (1969)& Intra-route & Micro Brute-force& A specialized constraint of $k$-permutation evaluating exactly $3! = 6$ permutations for clustered triples. & Ultra-fast micro-sequencing applied during post-crossover repair phases. \\ \midrule
{Chain Relocation} & Savelsbergh (1992) & Intra-route & Cascading Shift & Shifts a chain while temporarily permitting capacity/time violations, subsequently triggering forced secondary shifts. & Bypasses localized barriers via deterministic chain reactions; cascading evals risk exponential time. \\ \midrule

% --- INTER-ROUTE OPERATORS (LOCAL SEARCH) ---
{Swap* (Swap-Star)} & Vidal et al. (2012) & Inter-route & Amortized Swap & Swaps distinct nodes between routes while simultaneously evaluating optimal intra-route re-insertion. & $\mathcal{O}(n^2)$ evaluation achieved via sophisticated dynamic prefix/suffix arrays. A paradigm shift. \\ \midrule
{2-opt*} & Potvin \& Rousseau (1995)& Inter-route & Suffix Swap & Deletes one edge from Routes A/B, cross-connecting and permanently trading their downstream suffixes. & Highly effective for resolving scenarios where separate vehicles inefficiently cross paths. \\ \midrule
{3-opt*} & Potvin \& Rousseau (1995)& Inter-route & Complex Amalg. & Expands the cross-connection logic to three breaking edges across multiple active vehicles. & Rapidly executes large-scale vehicle amalgamation and inter-route segment transfers. \\ \midrule
{$k$-opt*} & Potvin \& Rousseau (1995)& Inter-route & Graph Shift & Extends macroscopic cross-edge reconnection across massive sets of vehicle architectures. & Generalization of star-operators; complexity scales extremely poorly without Potvin-style feasibility pruning. \\ \midrule
{Cross-Exchange} & Taillard et al. (1997) & Inter-route & String Trade & Swaps intact, contiguous sequences of customers between two completely separate vehicle paths. & Base $\mathcal{O}(n^4)$. Requires intelligent bounds or granular restriction to prune to $\mathcal{O}(n^2)$. \\ \midrule
{Imp. Cross-Exchange} & Bräysy (2003) & Inter-route & Amplified Trade & Extends base cross-exchange by forcing rapid, localized intra-route swaps directly upon traded segments. & Intensifies the topological perturbation to settle newly exchanged strings instantly. \\ \midrule
{$\lambda$-Interchange} & Osman (1993) & Inter-route & Bounded Swap & Methodically shifts or trades subsets of nodes (bounded rigidly by size $\le \lambda$) between two competing constraints. & Scales $\mathcal{O}(n^{2\lambda})$. Heavily reliant on dynamic programming piecewise linear penalty tracking. \\ \midrule
{Ejection Chain} & Glover (1991) & Inter-route & Feasib. Cascade & Initiates multi-route ripple effect: ejecting Node X to Route B forces Node Y from B to C, etc. & Phenomenal mechanism for bypassing severely bottlenecked constraint spaces. \\ \midrule
{Cyclic Transfer} & Thompson \& Psaraftis (1993)& Inter-route & Closed-loop & Enforces the multi-route ejection chain cascade to mathematically terminate in a closed loop. & Evaluates massive multi-vehicle load balances as a single differential delta. \\ \midrule
{(2,0)-Exchange} & Osman (1993) & Inter-route & Unbalanced Shift & Unreciprocated mathematical transfer of exactly two load elements strictly from Route A to Route B. & Core subset of the $\lambda$-interchange framework. Vital for shifting load away from max-capacity routes. \\ \midrule
{(2,1)-Exchange} & Osman (1993) & Inter-route & Imbalanced Swap & Transfers two nodes to a receiving route, but forcefully demands exactly one node transferred in return. & Navigates asymmetrical load balancing constraints dynamically. \\ \midrule
{$(k,0)$-Exchange} & Osman (1993) & Inter-route & Multi-node Purge & Pure shift of $k$ arbitrary nodes from one localized vehicle graph structure strictly to another. & The primary inter-route computational logic utilized to completely empty underutilized vehicles. \\ \midrule
{$(k,h)$-Exchange} & Osman (1993) & Inter-route & General Exchange & Fully generalized formulation transferring $k$ outgoing nodes in exchange for $h$ incoming nodes. & Navigates extreme edges of capacity constraints but computationally intensive if $k, h$ are unbound. \\ \midrule
{Split-Delivery Swap} & Dror \& Trudeau (1989) & Inter-route & Demand Partition & Explicitly breaks a single customer's demand across two or more vehicles dynamically. & Radically alters topology; breaks the foundational 1-node-1-route routing assumption (SDVRP). \\ \midrule

% --- PERTURBATION OPERATORS ---
{Kick (Ruin \& Recr.)} & Schrimpf et al. (2000) & Perturb & Large Shock & Massively executes destroy and repair on up to 30-40\% of the entire routing architecture. & A hard-reset mathematical mechanism intended to forcefully eject the search from basins. \\ \midrule
{Random Walk} & Lourenço et al. (2003) & Perturb & Local Degradation& Rapid series of random localized swap or shift variations that deliberately, slightly degrade the objective. & Small-scale escape velocity that breaks immediate local bounds in Iterated Local Search. \\ \midrule
{Double-Bridge} & Martin et al. (1991) & Perturb & 4-opt Cross & Removes four non-adjacent edges and reconnects isolated sub-tours in a non-sequential, crossed pattern. & Highly powerful trajectory tunnel; physically bypasses standard geometric bounds. \\ \midrule
{Elite-guided Trans.} & Vidal et al. (2012) & Perturb & Targeted Shock & Maps spatial vectors to a stored elite archive, using matrices to steer perturbation toward high-quality regions. & Resolves the computational waste of blind random jumps by genetically guiding the shock trajectory. \\ \midrule
{Evolutionary Perturb.} & Moscato (1989) & Perturb & Local Micro-GA & Isolates a sub-cluster of routes and rapidly executes crossover/mutation purely within that localized frame. & Merges intense localized trajectory search with the structural diversification of MAs. \\ \midrule
{Mutation (Inversion)} & Banzhaf (1990) & Perturb & Sequence Shift & Displaces or reverses a sub-string completely outside the context of classical 2-opt/3-opt cost tracking. & Unbounded structural shock utilized in classical GAs before local-search hybridizations. \\ \midrule

% --- NEURAL / ML OPERATORS (MODERN ADDITIONS) ---
{Neural Ruin (NLNS)} & Hottung \& Tierney (2020) & ML/Neural & Prob. Masking & Uses Attention to output softmax probabilities over nodes, aggressively removing those with low confidence. & Replaces hard-coded Shaw metrics with continuous latent representations of constraints. \\ \midrule
{Heatmap Crossover} & Joshi et al. (2019) & ML/Neural & Prob. Matrix & Uses Supervised Learning (GCNs) to generate heatmaps, restricting crossover only to probable edges. & Drastically cuts $\mathcal{O}(n^2)$ operations down to $\mathcal{O}(k)$ by preventing doomed edge evaluations. \\ \midrule
{NeuRewriter (RL)} & Chen et al. (2019) & ML/Neural & MDP Ejection & Replaces deterministic cyclic transfers with an RL agent actively learning the optimal node to eject. & Adapts dynamically to constraint bottlenecks that traditional programming struggles with. \\ \midrule
{DPDP Edge Pruning} & Kool et al. (2022) & ML/Neural & Search Pruning & Uses a GNN to score edges, then strictly limits Dynamic Programming or local search to the top-$k$ percentile. & Reduces exponential $\mathcal{O}(2^n)$ routing subsets to computationally tractable sparse graphs. \\ \midrule
{Active Search (EAS)} & Hottung et al. (2021) & ML/Neural & Test-time Adapt. & Updates the neural network's weights via RL gradients *during* inference on a single specific instance. & Blurs the line between Deep Learning inference and classical meta-heuristic trajectory search. \\ \midrule
{POMO Rollouts} & Kwon et al. (2020) & ML/Neural & Symm. Perturb. & Forces an autoregressive Attention Model to start constructive generation from $N$ different roots. & Highly effective exploitation of VRP graph symmetries without explicit search; $\mathcal{O}(N \times n)$. \\

\end{longtable}
\end{center}

\end{document}
